\section{备用资料}
\presenter{第 8 队}
\renewcommand{\pageindicator}{备用}
\source{原文表8，p.132。案例预测为论文报告值}
\begin{frame}[noframenumbering]{备用：共同参加活动的案例}
\renewcommand{\arraystretch}{1.5}
\par\noindent\begin{tabularx}{\linewidth}{@{}p{30mm}Y@{}}
\toprule
方法 & 回答\\\midrule
Flat Memory & 两人没有一起去过。\\
Graph Memory & 上个月，2023年6月。\\
StructMem & 去年，2022年8月。\\
参考答案 & \textcolor{accent}{2022年}\\
\bottomrule
\end{tabularx}\par
\takeaway{互动语境连接共同经历，来源日期帮助解析“去年”}
\aside{表中参考答案验证到年份，额外的月份没有得到单独验证。}
\note{现在回到开头的问题。平面记忆保留了两个人分别参加活动的事实，但没有恢复共同经历。图记忆也保留各自参加的边，不过这个案例中仍然没有正确连接两人。StructMem把互动条目与事实整合到同一叙述里，得到正确年份。要留意它的输出还带了月份，表中的标准答案只有2022，所以讲这个案例时，准确的结论应落在年份与共同参与关系上。}
\end{frame}

\source{原文图3(d)，p.126。纵轴截取75.4--77.05以显示细小差异}
\begin{frame}[noframenumbering]{备用：历史种子数量}
\begin{columns}[c,onlytextwidth]
\column{.65\textwidth}
\begin{tikzpicture}
\begin{axis}[width=8.8cm,height=5cm,xlabel={历史种子数量 $K$},ylabel={准确率（\%）},xmin=0,xmax=20,ymin=75.4,ymax=77.05,xtick={0,5,10,15,20},ytick={75.5,76,76.5,77},axis lines=left,grid=major,grid style={gray!15},tick label style={font=\scriptsize},label style={font=\small}]
\addplot[color=accent,mark=*,thick,nodes near coords,point meta=y,every node near coord/.append style={font=\scriptsize,anchor=south},/pgf/number format/fixed,/pgf/number format/precision=2] coordinates {(0,75.71)(5,76.43)(10,76.49)(15,76.82)(20,76.23)};
\end{axis}
\end{tikzpicture}
\column{.31\textwidth}
$K=15$时达到最高值。\par\vspace{4mm}
$K=0$仍整合当前缓存，只是不引入历史。\par\vspace{4mm}
更多历史不保证继续提升。
\end{columns}
\note{这张图测试构建阶段引入多少历史种子。15个时最高，20个时回落。0个依然会生成当前缓存的摘要，所以不能把它和上一页删除跨事件组件直接画等号。论文还发现，单纯扩大平面检索条数，在60条之后收益趋于平台。它支持增加合适范围的证据，也提醒我们关注上下文的组织与噪声。}
\end{frame}

\source{备用页：原文表4--5，p.130。表1口径复核详见阅读总结末页}
\begin{frame}[noframenumbering]{备用：更换评判模型后的结果}
\renewcommand{\arraystretch}{1.5}
\par\noindent\begin{tabularx}{\linewidth}{@{}Yrrr@{}}
\toprule 评判模型 & StructMem & Memobase & 分差\\\midrule
gpt-4o-mini & 76.82 & 75.78 & +1.04\\
Qwen2.5-32B & 77.01 & 76.36 & +0.65\\
DeepSeek-V3.2 & 79.35 & 79.29 & +0.06\\\bottomrule
\end{tabularx}\par
\vspace{4mm}
Fleiss' $\kappa=0.8341$，模型间判断较一致。\par\vspace{2mm}
分差随评判模型变化，DeepSeek下两者十分接近。
\aside{一致性指标衡量评判模型之间的共识，论文未报告系统分差的显著性检验。}
\note{如果被问到模型评判是否稳定，可以用这一页。三个评判模型下StructMem总分都略高，模型间一致性也较高。但DeepSeek下只领先0.06个百分点，因此应保留分差的实际量级。模型之间一致，并不直接表示它们与人工判断完全一致。阅读总结末尾还记录了少量表格加总差异，复现时需要用逐题输出澄清。}
\end{frame}

\source{备用页：原文 §3，附录 A.3、A.7，Limitations}
\begin{frame}[noframenumbering]{备用：几个方法细节}
\small
\textbf{为什么不用图也能称为结构化记忆？}\par
同源条目的时间绑定和跨事件摘要，已经提供了层次与关系组织。
\par\vspace{4mm}
\textbf{整合摘要会覆盖原来的记忆吗？}\par
事件条目与摘要共同保留，回答时同时检索。
\par\vspace{4mm}
\textbf{是否需要训练新模型？}\par
论文核心流程使用提示调用、向量检索与记忆组织，没有新的参数训练步骤。
\par\vspace{4mm}
\textbf{怎样处理相互冲突的记忆？}\par
问答提示优先较新信息，存储层的显式修订仍是待解决的问题。
\note{这一页用于回答方法层面的提问。结构化不只意味着实体关系图，也可以来自稳定的绑定关系和分层组织。StructMem保留条目与摘要，因此两者有互补作用。它不要求训练新的语言模型，适合用API完成小规模实验。最后要区分回答策略和存储更新，优先较新记忆有助于回答，但不能自动清除或修订所有旧摘要。}
\end{frame}
